% =========================================================
\chapter{Lecture 2 : Principles of computer security}
% =========================================================

\begin{rembox}[title=In a nutshell]
\textbf{What is this lecture about?} Nobody knows how to build a system without bugs, so
experience gave us \textbf{8 (+2) rules of thumb} (Saltzer \& Schroeder) to design
security mechanisms: keep it simple, refuse by default, check every access, do not rely on
secrecy of the design, never trust a single person, give only the rights needed, share as
little as possible, make security easy to use. \\
\textbf{Why is it useful?} These principles are a \textbf{toolbox to judge a design}: when
you look at a system (or an attack), you can ask ``which principle was violated?''. Most
real-world breaches (Heartbleed, MongoDB, S3 leaks, \ldots) are a principle that was
broken because of a bad tradeoff.
\end{rembox}

\begin{rembox}[title=Where do the principles come from?]
Saltzer \& Schroeder, \textit{The Protection of Information in Computer Systems} (1975),
abbreviated \textbf{[SS75]}: \textbf{8 (+2) principles} at the core of security
engineering practice. Experience has shown that they \textbf{guide} the design and
contribute to an implementation without security flaws. They are still used today
(e.g.\ to analyse the threat model of machine learning systems).
\end{rembox}

\begin{warnbox}[title=How to use the principles]
Always try to follow them. It may be impossible to follow all of them at once, but each must at
least be \textbf{considered}: if one is not followed, there must be a good reason, and the risks
of not following it must be understood.
\end{warnbox}

\begin{goatbox}
Nobody knows how to build a perfect barn. So old, wise farmers wrote down ten rules that make
barns break less often and less badly: keep it simple, keep the door closed unless told
otherwise, check everyone every time, do not count on the fox not knowing the plan, need two keys
for the big door, give each animal only the keys it needs, do not let animals share a trough,
make the locks easy enough that the farmer really uses them.
\end{goatbox}

\begin{exbox}[title=Heartbleed (CVE-2014-0160) --- why simplicity matters]
In OpenSSL (used for HTTPS), the \textit{Heartbeat} feature lets a client send a small
piece of data that the server echoes back. The attacker sends 1 byte but
\textbf{claims} to have sent 64\,000 bytes; the server, \textbf{without checking the
actual length}, replies with the byte plus 63\,999 bytes of its private memory.
\begin{itemize}
  \item \textbf{Feature complexity} --- Heartbeat was a non-essential feature: the
        added complexity created a \textbf{new attack surface}.
  \item \textbf{Codebase complexity} --- the bug was a single missing bounds check,
        hidden for two years because the huge codebase made auditing difficult.
\end{itemize}

The missing line was essentially \texttt{if (request\_size < actual\_size) \{ error; \}}. The
core cryptography of TLS was secure: the flaw was in an unsuspicious side feature. (Heartbleed
also started the fashion of vulnerabilities with a logo and a name, which raised public
awareness.)
\end{exbox}

\section{1 -- Economy of mechanism}

\begin{defnbox}[title=Economy of mechanism]
``Keep the [security mechanism / implementation] design as \textbf{simple and small}
as possible.'' (KISS principle)
\end{defnbox}

\begin{rembox}[title=Why?]
The mechanism needs to be \textbf{easy to audit and verify}. Operational testing is not
appropriate to evaluate security (penetration testing is nevertheless valuable).
\end{rembox}

\begin{rembox}[title={Why testing is not enough}]
Security is not about normal operation, nor about random faults: it is about what goes wrong when
the adversary chooses particular inputs, and one cannot test all inputs. Penetration testing
(trying to break the system) covers a bigger part of ``what could go wrong'' and is valuable, but
it can never give the \textbf{assurance} that nothing can go wrong. Only a mechanism small and
simple enough to be checked completely can.
\end{rembox}

\begin{defnbox}[title=Trusted Computing Base (TCB)]
\textbf{Every component} of the system (hardware / firmware / software) on which the
security policy relies. The TCB is \textbf{trusted} to operate correctly for the
security policy to hold.
\end{defnbox}

\begin{propbox}[title=Consequences of the TCB]
\begin{itemize}
  \item If something goes wrong \textbf{within} the TCB, \textbf{the security policy
        may be violated}; a failure \textbf{outside} the TCB should not matter for the
        policy.
  \item The TCB must be kept \textbf{small} to \textbf{ease verification} (economy of
        mechanism) and \textbf{diminish the attack surface}.
\end{itemize}
\end{propbox}

\begin{warnbox}[title=The verb ``to trust'']
The only proper use of ``to trust'' in security engineering is
\textbf{``X trusts Y will do Z''}.
\end{warnbox}

\begin{goatbox}
A big machine with a thousand gears hides a broken gear much better than a small one with five.
The TCB is the part of the barn you \emph{must} count on: the lock and the farmer who holds the
key. If anything outside it breaks (a window, a bucket), the hay is still safe; if the lock
breaks, it is not. So keep that part tiny, so you can check every screw.
\end{goatbox}

\section{2 -- Fail-safe defaults}

\begin{defnbox}[title=Fail-safe defaults]
``Base access decisions on \textbf{permission} rather than \textbf{exclusion}.'' [SS75]
\\
If something fails, be as secure as if it does not fail: errors / uncertainty should
err \textbf{on the side of the security policy}. Do \textbf{not} try to fix!
\end{defnbox}

\begin{rembox}[title={Why permission, and why not fix?}]
Permissions are given \emph{after thought}: if a permission exists, someone checked that the
action is allowed and safe. The space of non-permissions is \textbf{infinite}: one can never list
all harmful events in a blacklist. \\
After an error, the state of the system is unknown, so one cannot be sure that steps ``towards a
fix'' are safe: the safe move is to go back to a \textbf{known safe state} (deny, stop, roll
back).
\end{rembox}

\begin{propbox}[title={Whitelist, not blacklist}]
Use a \textbf{whitelist} (list what is allowed), not a blacklist (list what is
forbidden): a \textbf{lack of permission is easy to detect and solve}, whereas a
forgotten exclusion silently grants access.
\end{propbox}

\begin{exbox}[title=Fail-safe defaults]
\begin{itemize}
  \item Security door: if no permission, do not open.
  \item Form input: if no permission to write in $X$, do not write anywhere.
\end{itemize}
\end{exbox}

\begin{exbox}[title=MongoDB ransomware (2017) --- ease of use vs.\ security]
For a while, the default installation had \textbf{authentication turned off} and was
\textbf{bound to all network interfaces} (0.0.0.0) instead of localhost. Attackers
scanned the whole internet, connected without password, stole/deleted the data and left
a ransom note. Unsafe defaults + users who do not read the manual = no access control.
\end{exbox}

\begin{rembox}[title=Ransomware]
A ransomware \textbf{denies you access to your own data}, most commonly by
\textbf{encrypting your files}; the attacker then demands a \textbf{ransom payment} in
exchange for the key needed to decrypt them. Its prevalence exploded with the arrival of
cryptocurrencies.
\end{rembox}

\begin{exbox}[title={The first ransomware (1989)}]
20\,000 floppy disks labelled ``AIDS Information Introductory Diskette'' were mailed to attendees
of a WHO AIDS conference. After 90 boots, the malware encrypted (with simple, broken encryption)
all file names of the \texttt{C:} drive, hid directories, and asked for \$189 sent by post to a
P.O.\ box in Panama. Some victims panicked and wiped their drives: an Italian AIDS organisation
lost 10 years of research data. (A very similar attack later hit Elasticsearch databases left
open by default.)
\end{exbox}

\begin{goatbox}
The barn door is closed unless the farmer's list says ``this goat may enter''. If the list is
torn or the lock is jammed, the door stays \emph{closed}, never open. Writing ``wolves may not
enter'' is useless: tomorrow a bear comes, and he is not on the list of forbidden animals.
MongoDB shipped barns with the door wide open, expecting farmers to read the manual and close it.
They did not.
\end{goatbox}

\section{3 -- Complete mediation}

\begin{defnbox}[title=Complete mediation]
``\textbf{Every access to every object} must be checked for authority.'' [SS75]
\end{defnbox}

\begin{defnbox}[title=Reference monitor]
The component that \textbf{mediates ALL actions from subjects on objects} and ensures
they are according to the \textbf{policy}. It takes as input the request
(subject $\to$ object) and the policy, and writes to an \textbf{audit log}.
\end{defnbox}

\begin{warnbox}[title=Difficult to implement]
\begin{itemize}
  \item \textbf{Performance} --- checking everything is slow.
  \item \textbf{Time to check vs.\ time to use} --- the permission may change between
        the check and the use.
  \item \textbf{Modern distributed systems} --- you can only check what you see!
\end{itemize}
\end{warnbox}

\begin{rembox}[title=Why each difficulty]
\textbf{Performance}: think of how many memory accesses per second a computer makes, each one
checked. \textbf{Time of check vs.\ time of use}: between the check and the action the state may
change, so the check may be stale. \textbf{Distribution}: subjects and objects may be on
different machines and networks. A central monitor cannot see every access; a distributed one is
complex and hard to prove correct.
\end{rembox}

\begin{exbox}[title=Missing mediation --- the ``Gold Member'' cookie]
A streaming platform checks the subscription in the database \textbf{only once at
login}, then stores \texttt{\{'user': 'Alice', 'status': 'paid'\}} in a session cookie
and only looks at the cookie afterwards. Alice subscribes, logs in, unsubscribes (the DB
is correctly updated to ``unpaid''), and can still watch premium content until her
session expires: \textbf{violation of complete mediation}.
\end{exbox}

\begin{goatbox}
The guard must check the badge of every goat at every door, every time --- not only at the
gate in the morning. Otherwise a goat that was allowed in this morning but was kicked out at noon
keeps walking around all afternoon on its old ``I'm allowed'' sticker (the ``Gold Member''
cookie).
\end{goatbox}

\section{4 -- Open design}

\begin{defnbox}[title=Open design]
``The design should \textbf{not be secret}.'' [SS75] \\
When you design, \textbf{algorithms are public}; only \textbf{key elements} are kept
secret:
\begin{itemize}
  \item \textbf{Crypto} --- only keep the key secret;
  \item \textbf{Authentication} --- only keep the password secret;
  \item \textbf{Obfuscation} --- only keep the used noise secret.
\end{itemize}
\end{defnbox}

\begin{rembox}[title=Historical formulations]
\begin{itemize}
  \item \textbf{Kerckhoffs} (1883): the design of a system should not require secrecy.
        Every secret is a potential \textbf{failure point}: secrecy is a prime cause of weakness,
        openness means fewer points of failure.
  \item \textbf{Shannon} (1949): ``the enemy knows the system'' --- design systems
        assuming the enemy immediately gains full familiarity with them. Assuming the adversary
        does not know the system gives the adversary an advantage: he is \textbf{outside your
        threat model}.
  \item \textbf{Baran} (1964): the paradox of secrecy about secrecy --- without
        widespread scrutiny, weaknesses are more likely to be overlooked. Designing in secret
        prevents the conversations that find weaknesses.
\end{itemize}
\end{rembox}

\begin{propbox}[title=Why open design?]
\begin{itemize}
  \item Open design results in \textbf{better \& easier auditing}. \textbf{Linus' law}
        (Raymond, 1997): ``given enough eyeballs, all bugs are shallow'' --- more of a
        mantra: it needs many eyeballs, and defect density is high.
  \item \textbf{Secrecy is unrealistic}: assuming the design stays secret is a way to
        build a \textbf{bad threat model}.
  \item Key principle behind the academic discipline devoted to understanding computer
        security.
\end{itemize}
Famous failure of a closed design: \textbf{GSM encryption}.
\end{propbox}

\begin{exbox}[title={Closed designs that failed: GSM and DVD}]
\textbf{GSM encryption (1987)}: A5/1, with a 54-bit key (short on purpose, so that secret
services could brute-force it). It was reverse engineered and many serious attacks followed;
today A5/1 is decrypted in real time without the key (see also Lecture 5). \\
\textbf{DVD encryption (CSS, 1996)}: a proprietary 40-bit algorithm (US export law forbade longer
keys), given only under a non-disclosure licence. Every licensed player had a decryption key, and
each DVD's content key was encrypted for all players. One careless player maker did not protect
its key; hackers extracted it and could read (and copy) every DVD. Many secrets: one leak broke
the whole system. \\
\textbf{Defect density}: industry studies cite 15 to 50 bugs per 1000 lines of code, about 5\% of
them security vulnerabilities, i.e.\ roughly 1 per 1000 lines --- which is why ``enough
eyeballs'' is more of a mantra.
\end{exbox}

\begin{exbox}[title=Security through obscurity does not work]
\begin{itemize}
  \item \textbf{Hidden instance (2017)}: data of 198 million US voters stored on a
        public server with no authentication; the only ``protection'' was a
        ``semi-secret'' URL. A researcher found it.
  \item \textbf{Addresses are not safe}: the whole IPv4 space ($2^{32}\approx 4$
        billion addresses) could be scanned in about 3 minutes in 2013 --- picking a
        random port (0--65535) \textbf{is not hiding the service}.
\end{itemize}
\end{exbox}

\begin{goatbox}
Everybody may know how the lock is built; only the key is secret. A lock that is only safe
because nobody has seen it is a lock that breaks the day somebody looks at it --- and the fox
always looks. A hidden gate in the fence is still a gate: foxes walk along the whole fence in
three minutes.
\end{goatbox}

\section{5 -- Separation of privilege}

\begin{defnbox}[title=Privilege]
A \textbf{privilege} allows a user to perform an action on a computer system that may
have security consequences (e.g.\ create a file in a directory, access a device, write
to a socket for communicating over the Internet).
\end{defnbox}

\begin{defnbox}[title=Separation of privilege]
``No single accident, deception, or breach of trust is sufficient to compromise the
protected information.'' [SS75] \\
Requiring \textbf{multiple conditions} to execute an action improves security
(e.g.\ two keys to open a safe, two factors to authenticate).
\end{defnbox}

\begin{warnbox}[title=Problems]
\begin{itemize}
  \item \textbf{Availability?} --- all conditions are needed to use the service: what if you do
        not have your phone to log in to your bank?
  \item \textbf{Responsibility?} --- with several people involved, nobody in particular is
        responsible: it may lower security and makes liability harder to establish.
  \item \textbf{Complexity!} --- how are the conditions combined, when, does order matter?
        (Recall economy of mechanism.)
\end{itemize}
\end{warnbox}

\begin{exbox}[title=Payment system of a business]
\textbf{Goal}: the attacker must not steal money. \textbf{Threat model}: the attacker
can corrupt \emph{at most one} principal (no collusion). If a manager can both
\textit{create a vendor}, \textit{create an order} and \textit{approve} it, a single
corrupted manager registers a fake shell company and pays it. \\
\textbf{Solution: split roles} --- Procurement can \textbf{only} create/approve vendors;
Finance can \textbf{only} approve payments for vendors already in the system.
Reminiscent of \textbf{separation of powers} / ``checks and balances''.
\end{exbox}

\begin{exbox}[title=2AF]
A stolen password alone isn't enough.
\end{exbox}

\begin{goatbox}
The big safe needs two keys held by two different goats: one goat that goes bad, or is tricked,
or loses its key, is not enough to open it. The price: if one goat is on holiday, nobody opens
the safe.
\end{goatbox}

\section{6 -- Least privilege}

\begin{defnbox}[title=Least privilege]
``Every program and every user of the system should operate using the \textbf{least set
of privileges necessary} to complete the job.'' [SS75] \\
Rights are \textbf{added as needed} and \textbf{discarded after use}.
\end{defnbox}

\begin{propbox}[title=Least privilege in practice]
\begin{itemize}
  \item \textbf{Damage control} --- minimize high-privilege actions \& interactions
        (related to economy of mechanism: small TCB).
  \item \textbf{``Need-to-know'' principle} --- e.g.\ guest accounts at EPFL,
        \textbf{data minimization} principle (data protection).
\end{itemize}
\end{propbox}

\begin{rembox}[title={Which principle is it related to? Fail-safe defaults}]
A principal with few privileges cannot do much damage if it errs, misbehaves or is compromised:
even then, it cannot execute unauthorized actions, so the system stays in a \textbf{safe state}.
\end{rembox}

\begin{goatbox}
The goat that milks the cows gets the key of the milk room, not the key of the whole farm, and
gives it back when the milking is done. If that goat is tricked by the fox, the fox gets milk,
not the farm.
\end{goatbox}

\section{7 -- Least common mechanism}

\begin{defnbox}[title=Least common mechanism]
``Minimize the amount of mechanism \textbf{common to more than one user} and depended on
by all users.'' [SS75] \\
Every shared mechanism represents a \textbf{potential information path between users}
and must be designed with great care so that it does not unintentionally compromise
security.
\end{defnbox}

\begin{propbox}[title=Link with economy of mechanism]
\begin{itemize}
  \item \textbf{(Design)} interactions make it hard to validate the security design;
  \item \textbf{(Implementation)} interactions may lead to unintentional leaks of
        information --- \textbf{unintended channels}: use of \texttt{/tmp}, shared cache.
\end{itemize}
\end{propbox}

\begin{warnbox}[title={Mechanisms, not code}]
This principle refers to \textbf{mechanisms}, not to code: reusable code that has been
tested many times is \textbf{helpful} for security. The problem is the
\textbf{isolation of data}.
\end{warnbox}

\begin{exbox}[title=Two users sharing \texttt{/tmp}]
Alice and Bob both run programs on a multi-user server, which gives everyone the same
\texttt{/tmp}.
\begin{enumerate}
  \item Alice's program writes a temporary file \texttt{/tmp/report.tmp} containing her
        private salary data.
  \item Bob's program looks in \texttt{/tmp}, sees the file, and reads it (or overwrites
        it, or replaces it with a link to something else).
\end{enumerate}
Nobody wrote a bug that says ``let Bob read Alice's data'': \texttt{/tmp} is common to all
users, hence an \textbf{unintended information path} between them. \\
\textbf{Better design}: give each user a \textbf{private temp directory} (e.g.\
\texttt{/tmp/alice-xxxx/} with permissions only for Alice, or a per-user directory). Now
there is no shared channel to abuse.
\end{exbox}

\begin{exbox}[title={Exercise: find the vulnerability --- symlink attack in \texttt{/tmp}}]
\textbf{Principal 1}: a root program that writes a log to \texttt{/tmp/starting.log}.
\textbf{Principal 2}: the root logging subsystem. \textbf{Principal 3}: a user that runs non-root
programs, with write access to \texttt{/tmp} too. \textbf{Goal}: integrity of the files written
by principal 2. \textbf{Threat model}: the adversary can corrupt principal 3. \\
A \textbf{symbolic link} (symlink) is a ``shortcut'' file: when the OS accesses it, it is
automatically redirected to the real file it points to. \\
A root program writes \texttt{echo "Service is starting..." > /tmp/starting.log}. A
non-root user (corrupted by the adversary) also has write access to \texttt{/tmp} and
runs \texttt{ln -s /var/log/dmesg /tmp/starting.log}: the root program now overwrites a
system log, breaking \textbf{log integrity}. The common mechanism is the \texttt{/tmp}
directory, shared between low-privilege users and root processes. (Modern Linux
mitigates it in \texttt{/tmp}, but it still works e.g.\ in \texttt{/home/user}.)
\end{exbox}

\begin{goatbox}
If all the goats drink from one trough, a goat can leave a message in the water for another goat,
or poison it. Every shared thing (\texttt{/tmp}, a cache) is a secret tunnel between users. Give
each goat its own bucket. Sharing the \emph{recipe} of the bucket (reusing well-tested code) is
fine; sharing the \emph{water} is not.
\end{goatbox}

\section{8 -- Psychological acceptability}

\begin{defnbox}[title=Psychological acceptability]
``It is essential that the human interface be designed for \textbf{ease of use}, so that
users routinely and automatically apply the protection mechanisms correctly.'' [SS75]
\end{defnbox}

\begin{propbox}[title=What it implies]
\begin{itemize}
  \item \textbf{Hide the complexity} introduced by security mechanisms: they should not
        make the resource more difficult to access than if they were not present.
  \item The \textbf{mental model} of the (honest) users must match the security policy
        and security mechanisms.
  \item \textbf{Cultural acceptability} --- not all mechanisms are acceptable everywhere
        (face recognition in cultures that cover the face, register of everyone who
        sleeps in a dorm).
\end{itemize}
\end{propbox}

\begin{rembox}[title=Why it matters]
When mechanisms are complicated, users adopt \textbf{insecure practices} (the canonical example:
the password on a sticky note). If honest users do not understand the threat, they misuse the
mechanism; if it goes against social norms, they ignore it or circumvent it.
\end{rembox}

\begin{exbox}[title=Amazon S3 ``Any Authenticated AWS User'']
Psychological acceptability is \textbf{not just for end users}: system administrators
are also users of a configuration UI. Admins thought the option meant ``any
authenticated user \emph{in my account}'' whereas it meant ``anyone with \emph{any} AWS
account'' (free to create) $\Rightarrow$ several leaks from 2017. In 2023 AWS added a
fail-safe default (block public access) --- \textbf{safer defaults have tradeoffs}
(users struggled to share buckets).
\end{exbox}

\begin{goatbox}
A lock so annoying that the farmer props the door open with a stone protects nothing. The
safest lock is the one people actually use correctly without thinking --- and that includes the
farm's administrators reading a confusing settings screen (``any authenticated AWS user'' meant
\emph{any goat on Earth}).
\end{goatbox}

\section{Extra principles from physical security}

\begin{warnbox}[title=Difficult to transpose to computer security!]
The last two principles come from physical security and are hard to apply to computer
systems.
\end{warnbox}

\begin{defnbox}[title=9 -- Work factor]
``Compare the \textbf{cost of circumventing} the mechanism with the \textbf{resources of
a potential attacker}.'' [SS75] \\
It helps \textbf{refining the threat model}, but cost is \textbf{difficult to
quantify} (cost of compromising insiders? of finding a bug? monetization?).
\end{defnbox}

\begin{rembox}[title=Why it transposes badly]
In the physical world, cost is visible (breaking a lock, climbing a wall) and the benefit is
limited to the people physically there. In computers, what is the cost of corrupting an employee,
or of finding a bug (time? expertise?), and what is the benefit, when an exploit found once can be
used by many people at almost no cost?
\end{rembox}

\begin{defnbox}[title=10 -- Compromise recording]
``\textbf{Reliably record} that a compromise of information has occurred [\ldots] in
place of more elaborate mechanisms that completely prevent loss.'' [SS75] \\
Keep \textbf{tamper-evident logs}: they may enable recovery (integrity).
\end{defnbox}

\begin{warnbox}[title=Logs are not magic]
\begin{itemize}
  \item \textbf{What if you cannot recover?} If the attack deleted keys or leaked secrets,
        knowing that it happened (and who did it) does not bring them back.
  \item \textbf{How to keep the log's integrity?} If the adversary can tamper with the log, the
        log is useless: secure logs need their own security mechanism.
  \item \textbf{Logs may be a vulnerability (privacy)}: a log showing a patient talking very often
        to an oncologist reveals their health even if the records are confidential.
  \item \textbf{Logging the log? (availability)}: the log must not be deletable or blocked.
\end{itemize}
\textbf{Logging is not a guarantee that the compromise is detected.}
\end{warnbox}

\begin{goatbox}
Work factor: make the fox's job cost more than the chickens are worth to him --- but in computers
it is hard to price the fox's job. Compromise recording: if you cannot stop every theft, at least
keep a diary that the thief cannot erase, so you know what was taken and can sometimes get it
back.
\end{goatbox}

\begin{exbox}[title={The principles still apply today: machine learning}]
\textbf{Least privilege}: let the model learn as little as possible, so information cannot be
extracted from it. \textbf{Least common mechanism}: get samples labelled from different origins.
\textbf{Psychological acceptability}: users must understand why a model classifies or
misclassifies an input. \textbf{Work factor}: the cost of an attack (e.g.\ the number of API
calls) matters for its relevance. \textbf{Compromise recording}: ideally, log every step inside
the algorithm.
\end{exbox}

\begin{keybox}[title=Summary]
``Since no one knows how to build a system without flaws, the alternative is to rely on
eight design principles, which tend to reduce both the number and the seriousness of any
flaws.'' [SS75]
\begin{itemize}
  \item \textbf{Economy of mechanism} --- keep it simple!
  \item \textbf{Fail-safe defaults} --- safe defaults; if there is a problem, your move
        should comply with the policy.
  \item \textbf{Complete mediation} --- verify \emph{every} action.
  \item \textbf{Open design} --- make the design (and implementation) available.
  \item \textbf{Separation of privilege} --- never rely on \emph{only one} entity or
        action.
  \item \textbf{Least privilege} --- only the privileges needed, only while needed.
  \item \textbf{Least common mechanism} --- minimize shared mechanisms (isolate data).
  \item \textbf{Psychological acceptability} --- security must be easy to use correctly.
\end{itemize}
Principles help identify safe and unsafe \textbf{patterns}: do \textbf{not} use them as a
blind checklist, but as tools to weigh design decisions. They are easy to violate ---
typically through \textbf{improperly evaluated tradeoffs}, e.g.\ a default configuration made
easier to use (no certificate needed to host a website), or a simple, minimal TCB vs.\ the
simplicity of writing the code.
\end{keybox}

\begin{recapbox}[title={Exam-ready recap --- Lecture 2: Principles}]
Saltzer \& Schroeder (1975): \textbf{8 (+2) design principles}, to consider always (break one only
with a good reason and known risks); a tool to weigh tradeoffs, not a blind checklist. For an
exam scenario: name the \textbf{violated principle} and the \textbf{fix}.
\begin{center}
\small
\begin{tabularx}{\linewidth}{@{}>{\raggedright\arraybackslash}p{0.25\linewidth} >{\raggedright\arraybackslash}X >{\raggedright\arraybackslash}p{0.27\linewidth}@{}}
\toprule
\textbf{Principle} & \textbf{In one line} & \textbf{Example / violation} \\
\midrule
1 Economy of mechanism & simple, small mechanisms; small \textbf{TCB} (what the policy relies on; ``X trusts Y will do Z''); testing gives no assurance & Heartbleed (side feature, missing bounds check) \\
2 Fail-safe defaults & decide on \textbf{permission} (whitelist), not exclusion; on error, fall back to a safe state, do not fix & MongoDB open by default \\
3 Complete mediation & check \textbf{every} access to \textbf{every} object (reference monitor + audit log); hard: speed, TOCTOU, distribution & ``Gold Member'' cookie checked once \\
4 Open design & only keys / passwords are secret (Kerckhoffs, Shannon, Baran); obscurity fails & GSM A5/1, DVD CSS, hidden S3 URL, IPv4 scanned in 3 min \\
5 Separation of privilege & no single accident / deception / breach suffices; several conditions (cost: availability, responsibility, complexity) & 2FA, split Procurement / Finance \\
6 Least privilege & minimal rights, added when needed, dropped after; need-to-know; related to fail-safe & guest accounts, data minimization \\
7 Least common mechanism & every shared mechanism is an information path; isolate \emph{data} (shared code is fine) & \texttt{/tmp}, symlink attack, shared cache \\
8 Psychological acceptability & easy to use correctly; mental model matches the policy; culturally acceptable; admins are users too & S3 ``Any Authenticated AWS User'' \\
9 Work factor & cost of the attack vs.\ attacker's resources; hard to quantify in IT & \\
10 Compromise recording & tamper-evident logs; not magic (recovery, integrity, privacy, availability) & \\
\bottomrule
\end{tabularx}
\end{center}
Principles often pull against each other (a central reference monitor: economy of mechanism and
complete mediation vs.\ least common mechanism, Lecture 3).

\textbf{How it builds on what came before}: Lecture 1 said \emph{what} a secure system is (a
policy that holds under a threat model, shown by a security argument). This lecture says
\emph{how to design} the mechanisms so that such an argument is possible: small TCB, safe
defaults, every access checked, no reliance on secrecy of the design. The next lectures apply
these principles to the first concrete mechanism, \textbf{access control}: the reference monitor
\emph{is} complete mediation, negative permissions first \emph{is} fail-safe, roles and need-to-know
\emph{are} least privilege, covert channels \emph{are} least common mechanism.
\end{recapbox}
